\section{Casi di test}
% Ogni caso di test è descritto almeno da: numero identificativo, descrizione
% (anche la funzionalità testata), precondizioni, valori di input, output attesi
% e postcondizioni attese. All'esecuzione si aggiungono output e postcondizioni
% riscontrati ed esito (positivo se almeno un'asserzione non è rispettata,
% cioè il test ha rilevato un malfunzionamento; negativo altrimenti).

% Modello di tabella da riutilizzare per ogni gruppo di casi di test.
\begin{longtblr}[entry=none,label=none]{width=\linewidth,
             colspec={Q[l,wd=1.1cm]X[l]X[l]X[l]X[l]X[l]Q[l,wd=1.1cm]Q[l,wd=1.3cm]},
             rowhead=1, rowsep=2pt, cells={font=\scriptsize},
             row{1}={bg=arancio}, hline{2-Z}={wd=0.5pt, fg=arancio}}
TC & Precond. & Input & Output attesi & Output riscontrato & Postcond. & Esito & Priorità \\
TC-01 & % TODO
 & & & & & & \\
\end{longtblr}
